\section{Aula 4}

Nessa aula estávamos interessado em modelos em que podemos
estimar todas as probabilidades críticas explícitamente.
Exemplos destes modelos são as árvores $d$-árias infinitas.

\begin{definition}
	A árvore $d$-ária infinita $\cT_d$ é uma árvore com raiz $o$ em que
	todos os vértices tem $d$ filhos.
\end{definition}

O primeiro teorema da aula nos diz que sabemos $p_c$ e $\Theta(p_c)$
neste modelo.
Para essa seção em particular, seja $p_c^d$ a probabilidade crítica
em bond percolation de $\cT_d$.
\begin{theorem}
	Temos que
	\begin{equation}
		p_c^d = \frac{1}{d}
	\end{equation}
	e vale continuidade $\Theta_{\cT_d}^o(p_c^d) = 0$.
\end{theorem}

A prova é bem legal e usa uma exploração na árvore. Podemos fazer um
BFS ou DFS, tanto faz. Começamos
com $E = \varnothing$ e $L = \{o\}$, a cada passo do algoritmo,
escolhemos um vértice $v$ em $L$ e verificamos
se $v$ está conectado aos seus filhos (cada um com probabilidade $p$),
se um filho $w$ está conectado, adicionamos ele a $L$.
Depois de verificar todos os filhos de $v$, removemos $v$ de $L$ e o
adicionamos para $E$.

É fácil ver que esse processo explora a componente conexa da raiz.
Temos que verificar se o processo morre, i.e $L$ fica vazio em algum
momento. Basta notar que ele é equivalente a uma random walk,
começando do $1$ em que cada passo adiciona $Bin(p,d) - 1$ independentemente.

Segue que se $pd < 1$, o processo morre quase certamente e se $pd >
	1$ o processo sobrevive com probabilidade positiva. Uma forma de ver isso
é com as seguintes contas:

Definindo $S_n = 1 + \sum_{i \leq n} Bin(p,d) - 1$, vemos que se $pd
	- 1 < 0$ e usando Chernoff, segue que $\Pr(Sn \geq 0) < e^{-cn}$ e temos
morte q.c. Se $pd - 1 > 0$, então, por Chernoff de novo, existe $n_0$
tal que $\sum_{n \geq n_0} \Pr(S_n \leq 0) < Ke^{-cn_0} <1$.
Com probabilidade positiva, só crescemos até $n_0$ e nunca mais
morremos novamente.

O caso crítico é bem interessante e o Augusto só fez para $d = 2$.
Vamos ver que o processo morre, mas uma conta simples implica já que o número
de elementos na componente conexa tem esperança infinita. Seja $C_o$
o componente no final da exploração. Provaremos que
\[
	\frac{c}{\sqrt{n}} \leq \Pr(|C_o| \geq n) \leq C\sqrt{\frac{\log n}{n}}.
\]
Note que, durante a exploração, $|C_o| = S_n$ é uma martingale. Seja
$R = \inf \{t \in \N: S_t = 0 \lor S_t = \sqrt{n}\}$. Como a
sequência $S_{R_t} + 1$ é uniformemente
limitada, segue pelo teorema da parada que $1 = \Ex[S_m] = \Ex[1 +
		S_R] = \sqrt{n}\Pr(S_R = \sqrt{n})$, logo $\Pr(S_R = \sqrt{n}) = 1/\sqrt{n}$.
Agora temos que
\begin{align}
	\Pr(|C_o| \geq n) & \geq \Pr(S_R = \sqrt{n} \land 0 \not \in
	\{S_{R+1}, \dots S_{n}\})                                                                 \\
	                  & = \Pr(S_R \sqrt{n}) \Pr(\max_{i < n} S_i < \sqrt{n}) \geq \frac{c}{n}
\end{align}
onde a usamos o teorema de Donsker que nos diz
\[
	\Pr(\max_{i < n} S_i < \sqrt{n}) \to \Pr(N(0,\sigma) < \beta\sigma) = C.
\]
O upperbound é menos intuitivo. Separe o intervalo $[n]$ em $\log n$
subintervalos de tamanho $n/\log n$.
Seja defina $R' = \inf\{t \in \N : S_t = 0 \lor S_t = \sqrt{n/\log
		n}\}$, note que
\[
	\{|C_o| \geq n\} \subset \{S_{R'} = \sqrt{n/\log n}\} \cup \{\forall
	i \in [n], S_i \subset [\sqrt{n / \log n}]\}
\]
Como vimos antes a probabilidade do primeiro evento é $\sqrt{\log n /
		n}$. O segundo pode ser calculado
separando o intervalo $[n]$ em $\log n$ intervalos de tamanho $n/\log
	n$. Como espera-se que em cada um desses
o passeio pule $\sqrt{n/\log n}$ (em particular com probabilidade
constante $c < 1$ isso acontece), então
\[
	\Pr(|C_o| \geq n) \leq \sqrt{\frac{\log n}{n}} + [\Pr(\forall 1 \leq
	i \leq n/\log n : S_i < \sqrt{n/\log n})]^{\log n} \leq \sqrt{\frac{\log n}{n}} + c^{\log n}.
\]
O bound segue daí. Para obter o resultado formalmente, podemos multiplicamos os tamanhos dos intervalinhos
por alguma constante $c$, o mesmo para a altura do stopping time $R'$. Depois de otimizar conseguimos
o $C\sqrt{\frac{\log n}{n}}$.
